% --- Color Palette (Defined in Preamble) ---

\section{Introduction} \label{ch1}

Throughout this report, we will analyze the design and implementation of a hardware accelerator capable of applying Sobel filtering to monochromatic images of size 352$\times$288 pixels. The intensity of a pixel is represented by an 8-bit number in the integer range from 0 to 255, where 0 and 255 refer to black and white, respectively.

\subsection{Sobel Filtering}

The Sobel filter is a technique used to detect edges by computing the spatial gradient of an image; it detects sharp changes in brightness, thus highlighting boundaries between objects and backgrounds. It works by using a method called convolution: given a horizontal ($G_x$) and a vertical ($G_y$) kernel
$$G_x=\begin{bmatrix}
    -1 & 0 & 1 \\
    -2 & 0 & 2 \\
    -1 & 0 & 1 \\
\end{bmatrix}\quad G_y=\begin{bmatrix}
    1 & 2 & 1 \\
    0 & 0 & 0 \\
    -1 & -2 & -1 \\
\end{bmatrix}$$
The values of the derivatives $D_x$ and $D_y$ of a pixel $n$ are computed as convolution.
\begin{align*}
    D_x(n)=G_x * S(n)\\
    D_y(n)=G_y * S(n)
\end{align*}
To avoid the use of complex square root circuits, we approximate the length of the gradient with $L1$ magnitude.
$$|D(n)|=|D_x(n)|+|D_y(n)|$$
The final value $D(n)$ is then adapted to the $[0,255]$ pixel range.

\subsection{System Overview}

What is the setup?